\documentclass[a4paper,10pt]{article}

\usepackage[utf8]{inputenc}
\usepackage[T1]{fontenc}
\usepackage[margin=0.85in]{geometry}
\usepackage{enumitem}
\usepackage{titlesec}
\usepackage{graphicx}
\usepackage{parskip}
\usepackage{xcolor}
\usepackage[colorlinks=true, linkcolor=black, urlcolor=blue, citecolor=black]{hyperref}

\pagestyle{empty}
\setlength{\parindent}{0pt}
\renewcommand\labelitemi{--}
\renewcommand\labelitemii{$\circ$}
\titleformat{\section}{\large\bfseries}{}{0em}{}[\titlerule]

\begin{document}

%-------------------- HEADER --------------------
\begin{minipage}[c]{0.74\textwidth}
    {\Huge \textbf{Amirhossein Donyadidegan}}\\[0.25em]
    {\normalsize Jobtitel \textbar{} Stichwort 1 \textbar{} Stichwort 2 \textbar{} Stichwort 3}\\[0.5em]
    Karlsruhe, Baden-W\"urttemberg, Deutschland\\
    \href{mailto:donyadideganamir@gmail.com}{donyadideganamir@gmail.com}
    \quad | \quad +49 176 2477 0243\\
    GitHub: \href{https://github.com/AmirDonyadide}{AmirDonyadide}
    \quad | \quad
    LinkedIn: \href{https://www.linkedin.com/in/amirhossein-donyadidegan/}{amirhossein-donyadidegan}
\end{minipage}
\hfill
\begin{minipage}[c]{0.22\textwidth}
    \centering
    \includegraphics[width=1.05\linewidth]{photo.png}
\end{minipage}

%-------------------- PROFILE --------------------
\section*{Profil}
Berufseinsteiger im Bereich Geoinformatik mit abgeschlossenem MSc in Geoinformatik-Ingenieurwesen und praktischer Erfahrung in GIS, Fernerkundung, Python-Entwicklung, Datenanalyse und maschinellem Lernen. Meine Arbeit verbindet Geodatenwissenschaft mit datengetriebenen Methoden, einschließlich räumlicher Analysen, Dashboard-Entwicklung, Earth-Observation-Workflows und Machine-Learning-Anwendungen im geospatialen Kontext.

Meine Masterarbeit, \textit{Inferring Map Generalization Operations from User Prompts}, untersuchte, wie kartografische Generalisierungsprozesse mit KI und natürlicher Sprache verknüpft werden können, um Entscheidungsprozesse im Kartendesign zu unterstützen. Neben meiner akademischen Tätigkeit war ich an angewandten Forschungsprojekten am Karlsruher Institut für Technologie (KIT) beteiligt, wo ich Python-basierte Datenpipelines, geospatiale Analysen und interaktive Dashboards für industrielle Energie- und Mobilitätsanwendungen entwickelte sowie Web-/XR- und Visualisierungsprojekte unterstützte.

Ich suche Einstiegspositionen in den Bereichen Geoinformatik, GIS, Fernerkundung, geospatiale Datenanalyse, Data Science, maschinelles Lernen oder Python-basierte Entwicklung.

%-------------------- EDUCATION --------------------
\section*{Ausbildung}

\textbf{Politecnico di Milano}, Italien \hfill \textit{Sep. 2023 -- Mär. 2026}\\
MSc Geoinformatik-Ingenieurwesen, abgeschlossen \hfill Note: 102 / 110 ($\approx$ 1,5)

\textbf{Austauschprogramme:}
\begin{itemize}[leftmargin=*, itemsep=0.2em]
    \item \textbf{Universität Bonn}, Deutschland \hfill \textit{Okt. 2025 -- Mär. 2026}\\
    Erasmus+ Austauschstudent -- Geodäsie\\
    Masterarbeit: \textit{Inferring Map Generalization Operations from User Prompts}\\
    
    \item \textbf{Karlsruher Institut für Technologie (KIT)}, Deutschland \hfill \textit{Apr. 2025 -- Sep. 2025}\\
    Erasmus+ Austauschstudent -- Fernerkundung und Geoinformation
\end{itemize}

\textit{Relevante Kurse:} Geographic Information Systems, Machine Learning, Databases, Earth Observation, Geospatial Data Analysis, Geospatial Processing

\vspace{0.4em}

\textbf{Universität Teheran}, Iran \hfill \textit{Sep. 2018 -- Jul. 2022}\\
BSc Vermessungsingenieurwesen \hfill Note: 16.5 / 20 ($\approx$ 1,9)\\
Bachelorarbeit: \textit{Application of GIS and Big Data in Smart Cities}

\textit{Relevante Kurse:} Geographic Information Systems, Machine Learning, Databases, Earth Observation, Geospatial Data Analysis, Geospatial Processing

%-------------------- EXPERIENCE --------------------
\section*{Berufserfahrung}

\textbf{GIS- \& Datenanalyst (Hilfswissenschaftler) -- IIP, KIT} \hfill \textit{Jul. 2025 -- Mar. 2026}\\
\begin{itemize}[leftmargin=1.5em, noitemsep]
    \item Entwicklung von Python-basierten Datenpipelines zur Verarbeitung und Analyse von Energie- und Mobilitätsdaten
    \item Durchführung räumlicher Analysen zur Unterstützung von Energiebedarfs- und Dekarbonisierungsmodellen
    \item Entwicklung interaktiver Dashboards mit Kartenintegration für datengetriebene Anwendungen
    \item Kombination von Geodatenanalyse, Datenverarbeitung und Visualisierung in angewandten Forschungsprojekten
\end{itemize}

\vspace{0.4em}

\textbf{Web-Entwickler (Hilfswissenschaftler) -- KIT nova} \hfill \textit{Jul. 2025 -- Dez. 2025}\\
\begin{itemize}[leftmargin=1.5em, noitemsep]
    \item Mitarbeit an der Entwicklung und Verbesserung von Webanwendungen mit Fokus auf Funktionalität und Benutzerfreundlichkeit
    \item Unterstützung von VR/AR-Projekten für digitale und interaktive Anwendungen
    \item Beitrag zu geospatialen Visualisierungen und digitalen Kartenlösungen
\end{itemize}

\vspace{0.4em}

\textbf{Praktikant -- Naghsheh Gostaran Fartak Co.} \hfill \textit{Mai 2021 -- Sep. 2021}\\
\begin{itemize}[leftmargin=1.5em, noitemsep]
    \item Mitarbeit in Vermessung, Photogrammetrie und GIS-Projekten in einem praktischen Arbeitsumfeld
    \item Unterstützung bei Datenerfassung, räumlicher Analyse und technischer Umsetzung
    \item Anwendung von AutoCAD und GIS-Tools zur Erstellung und Verarbeitung von Geodaten
\end{itemize}

%-------------------- PROJECTS --------------------
\section*{Projekte}

\textbf{Inferring Map Generalization Operations from User Prompts} \hfill \href{https://github.com/AmirDonyadide/nl2map-generalization}{GitHub}\\
\begin{itemize}[leftmargin=1.5em, noitemsep]
    \item Entwicklung eines Machine-Learning-Workflows zur Verknüpfung von natürlicher Sprache mit kartographischen Generalisierungsoperationen
    \item Kombination von Text-Embeddings und geometrischen Features zur Vorhersage von Operatoren und Parametern
    \item Implementierung einer End-to-End-Pipeline inklusive Datenaufbereitung, Modelltraining und Evaluation
    \item Technologien: Python, Scikit-learn, GeoPandas, Shapely
\end{itemize}

\vspace{0.35em}

\textbf{LayerAlterator -- Raster Simulation Tool} \hfill \href{https://github.com/AmirDonyadide/LayerAlterator}{GitHub}\\
\begin{itemize}[leftmargin=1.5em, noitemsep]
    \item Entwicklung eines Python-Tools zur regelbasierten Modifikation von Rasterdaten mithilfe von Vektormasken
    \item Implementierung von Maskierungs- und prozentbasierten Transformationslogiken für räumliche Simulationen
    \item Sicherstellung von CRS-Konsistenz, Datenvalidierung und modularer Erweiterbarkeit
    \item Technologien: Python, Rasterio, GeoPandas, NumPy
\end{itemize}

\vspace{0.35em}

\textbf{DL4CVRS -- Vehicle Detection Model} \hfill \href{https://github.com/AmirDonyadide/DL4CVRS}{GitHub}\\
\begin{itemize}[leftmargin=1.5em, noitemsep]
    \item Implementierung und Vergleich eines Convolutional Neural Networks und eines vortrainierten ResNet18-Modells zur Bildklassifikation
    \item Aufbau einer vollständigen Trainings- und Evaluationspipeline mit Normalisierung, Metriken und Fehleranalyse
    \item Analyse von Fehlklassifikationen zur Bewertung und Verbesserung der Modellleistung
    \item Technologien: Python, PyTorch, TensorBoard, Scikit-learn
\end{itemize}

\vspace{0.35em}

\textbf{AI-Based Landslide Susceptibility Mapping Website} \hfill \href{https://github.com/AmirDonyadide/GIS-Course-Polimi-2024}{GitHub}\\
\begin{itemize}[leftmargin=1.5em, noitemsep]
    \item Entwicklung eines geospatialen Machine-Learning- und GIS-Workflows zur Analyse von Hangrutschungsanfälligkeit
    \item Integration von Fernerkundungs-, Terrain- und Infrastrukturdaten für Modellierung und Expositionsbewertung
    \item Umsetzung einer WebGIS-Anwendung zur interaktiven Visualisierung der Ergebnisse
    \item Technologien: GIS, Fernerkundung, JavaScript, WebGIS
\end{itemize}

\vspace{0.35em}

\textbf{LandsatToolkit -- Satellite Image Processing Library} \hfill \href{https://github.com/AmirDonyadide/LandsatToolkit}{GitHub}\\
\begin{itemize}[leftmargin=1.5em, noitemsep]
    \item Entwicklung einer Python-Bibliothek zur Verarbeitung und Analyse von Landsat-7/8/9-Satellitendaten
    \item Implementierung von Funktionen für Metadatenextraktion, Bandverarbeitung, Reprojektion und Indexberechnung
    \item Unterstützung wiederverwendbarer Earth-Observation-Workflows für Geodatenanalyse und Umweltmonitoring
    \item Technologien: Python, Rasterio, NumPy
\end{itemize}

\vspace{0.35em}

\textbf{SE4G -- Geospatial Data Visualization Dashboard} \hfill \href{https://github.com/AmirDonyadide/SE4G}{GitHub}\\
\begin{itemize}[leftmargin=1.5em, noitemsep]
    \item Entwicklung eines interaktiven Dashboards zur Visualisierung geospatialer Daten mit Karten und Diagrammen
    \item Integration von APIs und dynamischen Benutzerinteraktionen zur explorativen Datenanalyse
    \item Kombination von Backend- und Frontend-Komponenten in einer datengetriebenen Anwendung
    \item Technologien: Python, Dash, Plotly, Flask
\end{itemize}

\vspace{0.35em}

\textbf{PoliYoga -- Responsive Web Application} \hfill \href{https://github.com/moeinp70/HYP_Yoga}{GitHub}\\
\begin{itemize}[leftmargin=1.5em, noitemsep]
    \item Mitarbeit an der Konzeption und Entwicklung einer responsiven Webplattform für eine Yoga-Akademie
    \item Umsetzung von Frontend-, UX- und datenbanknahen Funktionen für Profile, Aktivitäten und Nutzerinteraktionen
    \item Beitrag zu einer benutzerzentrierten und modularen Webanwendung mit dynamischen Inhalten
    \item Technologien: Nuxt.js, Vue, TypeScript, Supabase, SQL, Figma
\end{itemize}

\vspace{0.35em}

\textbf{EOAdvanced -- Water Area Analysis in Eastern Venice} \hfill \href{https://github.com/AmirDonyadide/EOAdvanced}{GitHub}\\
\begin{itemize}[leftmargin=1.5em, noitemsep]
    \item Analyse zeitlicher Veränderungen von Wasserflächen in der Region Ost-Venedig auf Basis von Satellitendaten
    \item Berechnung und Auswertung von NDWI-basierten Indikatoren zur Untersuchung von Umweltveränderungen
    \item Erstellung von Karten und statistischen Visualisierungen für die Analyse von Trends zwischen 2021 und 2024
    \item Technologien: Python, Google Earth Engine, JavaScript, Geemap
\end{itemize}

\vspace{0.35em}

\textbf{Ingenieurvermessung -- Universität Teheran (Campusaufnahme)}\\
\begin{itemize}[leftmargin=1.5em, noitemsep]
    \item Durchführung einer topographischen Vermessung des Campus der Fakultät für Ingenieurwissenschaften
    \item Datenerfassung von Gebäuden und bebauten Flächen mit Totalstation
    \item Verarbeitung und Strukturierung der Messdaten zur Erstellung von Lageplänen
    \item Technologien: Totalstation, Vermessungstechnik
\end{itemize}

\vspace{0.35em}

\textbf{Nivellement -- Höhenvermessung (Universitätsgelände)}\\
\begin{itemize}[leftmargin=1.5em, noitemsep]
    \item Durchführung von Höhenmessungen zur Bestimmung von Höhenunterschieden im Gelände
    \item Anwendung klassischer Nivellement-Methoden und Fehlerkontrolle
    \item Auswertung und Dokumentation der Messergebnisse
    \item Technologien: Nivellementgeräte, Vermessung
\end{itemize}

\vspace{0.35em}

\textbf{Straßenplanung zwischen zwei Städten (AutoCAD Civil 3D)}\\
\begin{itemize}[leftmargin=1.5em, noitemsep]
    \item Entwurf einer Straßenverbindung unter Berücksichtigung nationaler Ingenieurvorschriften
    \item Planung von Trassierung, Kurvengeometrie und Längsprofil
    \item Berechnung von Erdbewegungen und Volumen (Cut \& Fill)
    \item Umsetzung des Projekts mit AutoCAD Civil 3D
    \item Technologien: AutoCAD Civil 3D, Straßenplanung
\end{itemize}

\vspace{0.35em}

\textbf{3D-Laserscanning -- Punktwolkenaufnahme eines Gebäudes}\\
\begin{itemize}[leftmargin=1.5em, noitemsep]
    \item Durchführung eines 3D-Laserscans eines kompletten Gebäudegeschosses
    \item Erfassung und Verarbeitung von Punktwolkendaten
    \item Unterstützung bei der Erstellung digitaler 3D-Modelle aus Scandaten
    \item Technologien: Laserscanner, Punktwolkenverarbeitung
\end{itemize}
%-------------------- TECHNICAL SKILLS --------------------
\section*{Technische Fähigkeiten}

\begin{itemize}[leftmargin=*, itemsep=0.3em]

    \item \textbf{Programmierung:} Python (OOP, modulare Entwicklung, Package-Entwicklung), C++, JavaScript, SQL, MATLAB

\item \textbf{Machine Learning \& Data Science:} Supervised Learning, Unsupervised Learning, Semi-Supervised Learning, Reinforcement Learning (Grundlagen), Klassifikation \& Regression, Binary \& Multi-Class Classification, Regression Modeling, Clustering, Anomaly Detection, Dimensionality Reduction, Feature Engineering, Spatial Feature Engineering, Geospatial Feature Extraction, Feature Scaling \& Normalisierung, Datenvorverarbeitung, Datenaugmentation, Modellvalidierung (Train/Test, Cross-Validation), Hyperparameter-Optimierung, Overfitting \& Regularisierung, Ensemble Learning, Performance-Metriken (Accuracy, Precision, Recall, F1-Score, Confusion Matrix), Modellinterpretation, Feature Importance Analysis, Explainable AI (XAI) (Grundlagen), Transfer Learning, Representation Learning, Similarity Learning, Embedding-Modelle \& Feature-Repräsentationen, Vector Embeddings, Text Embeddings, NLP-basierte Feature-Repräsentation, Multimodale Datenverarbeitung, Zeitreihenanalyse (Grundlagen), End-to-End ML Pipelines, Reproduzierbare ML-Workflows, Experiment Tracking, Pipeline-Automatisierung, AI-assisted Data Analysis, Modellbereitstellung (Grundlagen), Prompt Engineering (Grundlagen), Retrieval-Augmented Concepts (Grundlagen), Scikit-learn, TensorFlow, PyTorch, PCA, UMAP, t-SNE, K-Means, Random Forest, Gradient Boosting, Support Vector Machines (SVM), k-Nearest Neighbors (kNN), Decision Trees,

Neuronale Netzwerke: Multi-Layer Perceptron (MLP), Convolutional Neural Networks (CNN), Deep Learning Workflows, Transfer Learning mit vortrainierten Modellen (z.\,B. ResNet), Grundlagen von Graph Neural Networks (GNN), Grundlagen von Recurrent Neural Networks (RNN), Autoencoder (Grundlagen), Large Language Models (LLMs), Prompt-basierte Modellinteraktion

\item \textbf{Geospatial AI \& Remote Sensing ML:} Geospatial Machine Learning, Remote Sensing ML, Earth Observation Analytics, Raster- und Vektordatenanalyse, räumliche Modellierung, Satellitendatenanalyse, räumliche Mustererkennung, GIS-basierte ML-Workflows, Geodatenintegration für ML-Anwendungen

    \item \textbf{Datenverarbeitung \& Data Engineering:} Datenpipelines, Pipeline-Design, ETL-Workflows (Extract, Transform, Load), Datenintegration, Datenbereinigung und -vorverarbeitung, Data Transformation Workflows, Pandas, NumPy, GeoPandas, xarray, Jupyter Notebooks

    \item \textbf{GIS \& Fernerkundung:} QGIS, ArcGIS, Google Earth Engine, Earth Observation (EO), Satellitendatenanalyse, räumliche Datenmodellierung, Geodatenintegration, Spatial Feature Extraction, Raster- und Vektordatenverarbeitung, CRS-Transformationen, Kartennetze und Projektionen, räumliche Operationen (Buffer, Overlay, Spatial Join), GeoJSON, Shapefiles, LiDAR-Grundlagen

    \item \textbf{Visualisierung \& Dashboards:} Plotly, Dash, Matplotlib, interaktive Dashboards, kartographische Visualisierung

    \item \textbf{Webentwicklung \& APIs:} HTML, CSS, JavaScript, Nuxt.js, Flask, REST APIs, API-Design, API-Integration, Datenabruf über APIs, Backend--Frontend-Integration, Client--Server-Architektur, JSON-basierte Datenverarbeitung

    \item \textbf{Datenbanken:} PostgreSQL, PostGIS, SQLite, Oracle, Supabase

    \item \textbf{DevOps \& Software Engineering:} Softwarearchitektur (Grundlagen), modulare Systementwicklung, Git, GitHub, GitLab, Git-Workflows (Branching), Linux/Bash, Docker, CI/CD (GitHub Actions), Debugging, Unit Testing (pytest), virtuelle Umgebungen (venv, conda), Code-Dokumentation

    \item \textbf{Cloud \& Automatisierung:} Google Cloud, Google Sheets API, Google Drive API, Skripting, Automatisierung

    \item \textbf{Engineering \& Design Tools:} AutoCAD, Civil3D, Figma

\end{itemize}
%-------------------- LANGUAGES --------------------
\section*{Sprachen}

Englisch (fließend), Deutsch (Mittelstufe), Italienisch (Mittelstufe)

%-------------------- ADDITIONAL --------------------
\section*{Weitere Informationen}

\begin{itemize}[leftmargin=*, itemsep=0.25em]
    \item Führerschein Klasse B
\end{itemize}

\end{document}
